% preamble.tex — shared preamble for every chapter (Stacks-project style).
% Each chapter file starts with % preamble.tex — shared preamble for every chapter (Stacks-project style).
% Each chapter file starts with % preamble.tex — shared preamble for every chapter (Stacks-project style).
% Each chapter file starts with % preamble.tex — shared preamble for every chapter (Stacks-project style).
% Each chapter file starts with \input{preamble} and is a complete,
% independently compilable document.

\documentclass[a4paper,11pt]{amsart}

\usepackage{amsmath,amssymb,mathtools}
\usepackage{lmodern}
\usepackage[T1]{fontenc}

% Cross-chapter references.  `xr' (v6+, which absorbed the old xr-hyper)
% reads the referenced chapter's .aux file at \externaldocument time.  A
% missing .aux is only a warning ("No file ... LABELS NOT IMPORTED") —
% the chapter still compiles; the refs fill in once the other chapter
% has been built.
\usepackage{xr}

% Clickable references in the output PDF.  All links black on purpose:
% colored links render as washed-out gray on the e-ink panel.
\usepackage[colorlinks=true,allcolors=black]{hyperref}

% One \externaldocument line per chapter; the prefix is the chapter file
% name plus "-".  From any OTHER chapter you then write
% \ref{<chapter>-<label>}.  Add a line here each time you create a
% chapter.  RULE: never give a chapter a name that is a prefix of
% another chapter's name (e.g. `groups' and `groups-extra') — RefTeX
% resolves external labels by prefix matching and would pick the wrong
% file.
%\externaldocument[example-]{example}

% Theorem environments: one shared counter, numbered within sections.
% amsart already includes the amsthm code — do NOT \usepackage{amsthm}.
\theoremstyle{plain}
\newtheorem{theorem}{Theorem}[section]
\newtheorem{proposition}[theorem]{Proposition}
\newtheorem{lemma}[theorem]{Lemma}
\newtheorem{corollary}[theorem]{Corollary}
\theoremstyle{definition}
\newtheorem{definition}[theorem]{Definition}
\newtheorem{example}[theorem]{Example}
\newtheorem{exercise}[theorem]{Exercise}
\theoremstyle{remark}
\newtheorem{remark}[theorem]{Remark}

\numberwithin{equation}{section}

%%% Local Variables:
%%% TeX-master: t
%%% End:
 and is a complete,
% independently compilable document.

\documentclass[a4paper,11pt]{amsart}

\usepackage{amsmath,amssymb,mathtools}
\usepackage{lmodern}
\usepackage[T1]{fontenc}

% Cross-chapter references.  `xr' (v6+, which absorbed the old xr-hyper)
% reads the referenced chapter's .aux file at \externaldocument time.  A
% missing .aux is only a warning ("No file ... LABELS NOT IMPORTED") —
% the chapter still compiles; the refs fill in once the other chapter
% has been built.
\usepackage{xr}

% Clickable references in the output PDF.  All links black on purpose:
% colored links render as washed-out gray on the e-ink panel.
\usepackage[colorlinks=true,allcolors=black]{hyperref}

% One \externaldocument line per chapter; the prefix is the chapter file
% name plus "-".  From any OTHER chapter you then write
% \ref{<chapter>-<label>}.  Add a line here each time you create a
% chapter.  RULE: never give a chapter a name that is a prefix of
% another chapter's name (e.g. `groups' and `groups-extra') — RefTeX
% resolves external labels by prefix matching and would pick the wrong
% file.
%\externaldocument[example-]{example}

% Theorem environments: one shared counter, numbered within sections.
% amsart already includes the amsthm code — do NOT \usepackage{amsthm}.
\theoremstyle{plain}
\newtheorem{theorem}{Theorem}[section]
\newtheorem{proposition}[theorem]{Proposition}
\newtheorem{lemma}[theorem]{Lemma}
\newtheorem{corollary}[theorem]{Corollary}
\theoremstyle{definition}
\newtheorem{definition}[theorem]{Definition}
\newtheorem{example}[theorem]{Example}
\newtheorem{exercise}[theorem]{Exercise}
\theoremstyle{remark}
\newtheorem{remark}[theorem]{Remark}

\numberwithin{equation}{section}

%%% Local Variables:
%%% TeX-master: t
%%% End:
 and is a complete,
% independently compilable document.

\documentclass[a4paper,11pt]{amsart}

\usepackage{amsmath,amssymb,mathtools}
\usepackage{lmodern}
\usepackage[T1]{fontenc}

% Cross-chapter references.  `xr' (v6+, which absorbed the old xr-hyper)
% reads the referenced chapter's .aux file at \externaldocument time.  A
% missing .aux is only a warning ("No file ... LABELS NOT IMPORTED") —
% the chapter still compiles; the refs fill in once the other chapter
% has been built.
\usepackage{xr}

% Clickable references in the output PDF.  All links black on purpose:
% colored links render as washed-out gray on the e-ink panel.
\usepackage[colorlinks=true,allcolors=black]{hyperref}

% One \externaldocument line per chapter; the prefix is the chapter file
% name plus "-".  From any OTHER chapter you then write
% \ref{<chapter>-<label>}.  Add a line here each time you create a
% chapter.  RULE: never give a chapter a name that is a prefix of
% another chapter's name (e.g. `groups' and `groups-extra') — RefTeX
% resolves external labels by prefix matching and would pick the wrong
% file.
%\externaldocument[example-]{example}

% Theorem environments: one shared counter, numbered within sections.
% amsart already includes the amsthm code — do NOT \usepackage{amsthm}.
\theoremstyle{plain}
\newtheorem{theorem}{Theorem}[section]
\newtheorem{proposition}[theorem]{Proposition}
\newtheorem{lemma}[theorem]{Lemma}
\newtheorem{corollary}[theorem]{Corollary}
\theoremstyle{definition}
\newtheorem{definition}[theorem]{Definition}
\newtheorem{example}[theorem]{Example}
\newtheorem{exercise}[theorem]{Exercise}
\theoremstyle{remark}
\newtheorem{remark}[theorem]{Remark}

\numberwithin{equation}{section}

%%% Local Variables:
%%% TeX-master: t
%%% End:
 and is a complete,
% independently compilable document.

\documentclass[a4paper,11pt]{amsart}

\usepackage{amsmath,amssymb,mathtools}
\usepackage{lmodern}
\usepackage[T1]{fontenc}

% Cross-chapter references.  `xr' (v6+, which absorbed the old xr-hyper)
% reads the referenced chapter's .aux file at \externaldocument time.  A
% missing .aux is only a warning ("No file ... LABELS NOT IMPORTED") —
% the chapter still compiles; the refs fill in once the other chapter
% has been built.
\usepackage{xr}

% Clickable references in the output PDF.  All links black on purpose:
% colored links render as washed-out gray on the e-ink panel.
\usepackage[colorlinks=true,allcolors=black]{hyperref}

% One \externaldocument line per chapter; the prefix is the chapter file
% name plus "-".  From any OTHER chapter you then write
% \ref{<chapter>-<label>}.  Add a line here each time you create a
% chapter.  RULE: never give a chapter a name that is a prefix of
% another chapter's name (e.g. `groups' and `groups-extra') — RefTeX
% resolves external labels by prefix matching and would pick the wrong
% file.
%\externaldocument[example-]{example}

% Theorem environments: one shared counter, numbered within sections.
% amsart already includes the amsthm code — do NOT \usepackage{amsthm}.
\theoremstyle{plain}
\newtheorem{theorem}{Theorem}[section]
\newtheorem{proposition}[theorem]{Proposition}
\newtheorem{lemma}[theorem]{Lemma}
\newtheorem{corollary}[theorem]{Corollary}
\theoremstyle{definition}
\newtheorem{definition}[theorem]{Definition}
\newtheorem{example}[theorem]{Example}
\newtheorem{exercise}[theorem]{Exercise}
\theoremstyle{remark}
\newtheorem{remark}[theorem]{Remark}

\numberwithin{equation}{section}

%%% Local Variables:
%%% TeX-master: t
%%% End:
